\chapter{Comparing characterization techniques}
\label{ch:char-compare}

\paragraph{In this chapter.} Which technique answers which question (matrix on a landscape page), with probing depth, resolution and sample
requirements, plus the standard ``two techniques to distinguish\ldots'' answers.



\section{Probing depth, resolution, requirements}
\begin{center}
\small
\renewcommand{\arraystretch}{1.25}
\begin{tabularx}{\linewidth}{@{}lXXX@{}}
\toprule
Technique & Probe $\to$ signal & Information depth / resolution & Requirements \\ \midrule
UV-vis & light $\to$ absorbed light & whole sample (bulk average) & dispersion or film; plasmonic metal or semiconductor \\
Fluorescence & light $\to$ emitted light & whole sample & emitting material (QDs) \\
XPS & X-rays $\to$ photoelectrons & \textbf{top 5--10 nm}; mm-sized spot & UHV (or NAP); solid \\
XRD & X-rays $\to$ diffracted X-rays & bulk ($\mu$m); averages many crystals & crystalline, $\gtrsim$ 2--3 nm, $\gtrsim$ few \% \\
TEM / HAADF & electrons $\to$ transmitted electrons & $\sim$0.1 nm lateral; projection through sample & thin sample on grid; vacuum \\
EDX & electrons $\to$ X-rays & nm-scale maps & in TEM/SEM \\
EELS/EFTEM & electrons $\to$ energy loss & nm-scale maps, oxidation state & thin sample \\
AFM & tip force $\to$ deflection & \textbf{sub-nm vertical}, lateral limited by tip & flat support; any material \\
STM & tip current & atomic & conducting, clean surface \\
\bottomrule
\end{tabularx}
\end{center}

\begin{landscape}
\section{Question $\to$ technique matrix}
\textbf{\checkmark\checkmark} = the best/direct answer, \textbf{\checkmark} = possible/indirect, -- = no.
\begin{center}
\renewcommand{\arraystretch}{1.3}
\setlength{\tabcolsep}{4pt}
\resizebox{\linewidth}{!}{%
\begin{tabular}{@{}lcccccccc@{}}
\toprule
Question & UV-vis & XPS & XRD & TEM & HRTEM/FT & EDX/EELS & AFM & STM \\ \midrule
Particle size (distribution) & \checkmark & -- & \checkmark\ (crystallite) & \checkmark\checkmark & \checkmark & -- & \checkmark\checkmark\ (height) & \checkmark \\
Shape & \checkmark\ (rods) & -- & -- & \checkmark\checkmark & \checkmark\checkmark & -- & \checkmark & \checkmark \\
Crystal structure / phase & -- & -- & \checkmark\checkmark & -- & \checkmark\checkmark\ (local) & -- & -- & -- \\
Lattice parameter / alloy composition & \checkmark\ (band position) & -- & \checkmark\checkmark & -- & \checkmark & \checkmark & -- & -- \\
Exposed planes, epitaxy & -- & -- & -- & -- & \checkmark\checkmark & -- & -- & \checkmark\ (surface) \\
Elemental composition & -- & \checkmark\checkmark\ (surface) & \checkmark\ (phases) & -- & -- & \checkmark\checkmark\ (local) & -- & -- \\
Oxidation state / electronic interaction & \checkmark\ (plasmon on/off) & \checkmark\checkmark & \checkmark\ (phase) & -- & -- & \checkmark\ (EELS) & -- & -- \\
Where each element is (core/shell, alloy) & \checkmark & \checkmark\ (surface vs bulk) & -- & \checkmark\ (Z-contrast) & -- & \checkmark\checkmark\ (maps) & -- & -- \\
Inside or outside (tube, pore) & -- & \checkmark & -- & \checkmark\ (tomography) & -- & \checkmark & \checkmark\ (outside only) & -- \\
Capping/organic shell & -- & \checkmark\ (C 1s, S 2p) & -- & -- (invisible) & -- & -- & \checkmark\checkmark & -- \\
Representative of whole sample & \checkmark\checkmark & \checkmark & \checkmark\checkmark & -- & -- & -- & -- & -- \\
\bottomrule
\end{tabular}}
\end{center}
\end{landscape}

\section{Typical pairings asked in exams}
\begin{keybox}[: standard answers]
\begin{tabularx}{\linewidth}{@{}>{\raggedright\arraybackslash}p{5cm}X@{}}
\toprule
To decide\ldots & Use \\ \midrule
Mixture vs solid solution (CeO$_2$+ZrO$_2$ vs Ce$_{1-x}$Zr$_x$O$_2$) & XRD (one shifted set vs two sets); HRTEM d-spacings of
single particles; STEM-EDX/EELS maps (Ce and Zr in the same particle or not); Raman\extra. \\
Metal vs its oxide (Ag vs Ag$_2$O, Cu vs CuO) & XRD (different structure/lattice); XPS (chemical shift, O 1s, satellites,
Auger); UV-vis (plasmon only for metal). \\
Core vs shell composition & EELS/EDX maps; HAADF Z-contrast; HRTEM fringes of each region; XPS before/after shell removal;
UV-vis before/after. \\
Alloy vs core--shell vs mixture & UV-vis (one intermediate band vs two); EDX/EELS line-scan/maps; XPS shifts; XRD only if
$a$ differ enough. \\
Particles inside vs outside a tube/pore & TEM tomography; XPS (surface-sensitive: outside particles give stronger signal);
SEM/AFM (only outside visible); EDX maps. \\
Capping shell present? & AFM height vs TEM diameter; XPS (S 2p, C 1s); TGA\extra. \\
Particle size & TEM histogram (number); XRD Scherrer (crystallite); AFM height; UV-vis (indirect). \\
Exposed facets / epitaxy & HRTEM + FT (fringes parallel to edges; continuity of fringes across the interface). \\
\bottomrule
\end{tabularx}
\end{keybox}

\section{Summary: the unique advantage of each technique}
\begin{itemize}
  \item \textbf{UV-vis}: fast, cheap, whole-sample view of plasmonic metals (metal, size, shape, alloy, aggregation).
  \item \textbf{XPS}: the \emph{only} routine technique for \textbf{oxidation state and surface composition}.
  \item \textbf{XRD}: the reference for \textbf{crystal phase and lattice parameter}, plus crystallite size, averaged over the bulk.
  \item \textbf{TEM}: \textbf{direct image} of size, shape and location. HRTEM adds local structure, facets and epitaxy;
        EDX/EELS add local chemistry.
  \item \textbf{AFM}: true \textbf{heights} on any surface (including organic shells), in air or liquid.
  \item \textbf{STM}: \textbf{atomic resolution} and electronic structure of conducting surfaces, and atom manipulation.
\end{itemize}
